\runninghead{Scope and References}
\section{Appendix: Scope and Qualifications}
\label{sec:scope}

{\small
\textbf{Edition, formulary and occurrence.} The study treats the Mass
\latin{In festo omnium Sanctorum}, \latin{Die 1 novembris}, I classis, of the
\work{Missale Romanum}, editio typica, Typis Polyglottis Vaticanis 1962,
printed pp.~702--703, nos.~3974--3983, as said on Sunday 1~November 2026 in
the universal 1962 calendar and no particular calendar, with the commemoration
of the Twenty-third Sunday after Pentecost, pp.~416--417, nos.~1720, 1726 and
1728. The precedence of the feast over the Sunday, the commemoration and its
order, the colour, the Gloria, the Credo and the last Gospel rest on the
Missal's own general rubrics (RG 14, 16, 91, 108, 109~a, 111~a, 120~d; RGMR
431, 434~a, 436--437, 475, 480--481, 505, 509--510), as
\texttt{research/context.md} records. The Preface is not settled: the feast
has none of its own, and the general rubrics (RGMR 482, 483, 494~b, 498) can
be read for the Preface of the Most Holy Trinity or for the common Preface.
No other rubrical source searched for the question, as
\texttt{research/scope.md} lists them, settles it. The study names both,
asserts neither, gives neither text, and builds nothing on either. The \latin{Asperges}, the Ordinary and the Office of the feast are
outside the formulary and are not treated.

\textbf{Text control.} Every Latin form was read on page images of the CMAA
facsimile of the 1962 typical edition and collated word for word with the
Tours printing of 1922 (Mame) for the feast's ten texts and with the Ratisbon
printing of 1862 (Pustet) for the Sunday's three orations; the collation and
the comparison with the Clementine Vulgate are in
\texttt{propers/verified.md}. The rubrics of the 1962 book are described and
not reproduced. The history of the Offertory's wording rests on the 1570 text
layer read at the feast and at two Masses of martyrs, and on the text layers
of the 1604, 1862 and 1889 printings; when its \latin{alleluia} entered is not
established.

\textbf{English.} The scriptural English is the Douay--Rheims in Challoner's
revision at the canonical verses, with the Missal's departures named beside
each text. The English of the Introit antiphon and of the seven orations is
the Lasance \work{New Roman Missal} of 1945, whose rubrics are those before
1960. The Lasance English of the Lesson's and the Gospel's liturgical
openings and of the Offertory is quoted on the Lasance's attribution only; it
has no row in the repository's register of translations. None of this English
is an approved liturgical translation of the 1962 Missal. Patristic and later
witnesses are quoted in the nineteenth-century English of the editions named
in the References, or in their Latin or French, and otherwise described; no
English is composed for a liturgical text.

\textbf{Reception.} The interpretations rest on the loci named in the text
and recorded, with editions and bounds, in \texttt{research/scope.md} and
\texttt{research/interpretations.md}. The sweep is bounded by the corpora named
there and is not complete; St Jerome on Matthew, St Bernard's sermons for the
feast, St Gregory of Nyssa on the Beatitudes and the Greek commentators on the
psalms beyond Basil were not read. No commentary on the Lesson's chapter by St
Augustine was used, and he is not presented as interpreting it; his three
loci on the Lesson's number and crowd (\work{De doctrina christiana}
III.35.51, \work{Letter XCIII} \S30, \work{Sermon LXI} 1) are short
applications inside other arguments. Victorinus of Pettau's commentary on
the chapter is quoted from the Ante-Nicene Fathers text under its editors'
recension caution --- corrupted and interpolated manuscripts, a possible
later corrector, an uncertain author --- and St Cyprian's martyrs-only
reading of the crowd is kept as the reading Bede rejects; Victorinus and
Cyprian have no rows in the author-standing register, and neither carries an
interpretation. Blessed
Ildefonso Schuster's \work{The Sacramentary} V (English, 1930), Lucien
Fromage's continuation of Guéranger's \work{L'Année liturgique} (vol.~15,
1909) and Durandus's \work{Rationale} are the commentaries read on this Mass
as a Mass; Fromage and Durandus are ecclesiastical writers and carry no
interpretation. Rabanus Maurus is cited as a compiler of Augustine. Every
connection drawn between texts of the Mass belongs to the study's
interpretations unless a witness is named as making it, and the orations of
the commemorated Sunday are read only for what they say where they stand in
this Mass. Bellarmine is read in O'Sullivan's English, which numbers Psalm~33
one verse lower than the Vulgate; Basil in the Latin version printed in PG
29; Aquinas in the Venice (1745) and Parma (1863) printings, on their text
layers.

\textbf{Chronology.} Every biblical date in the Scriptural Date and Location
appendix is generated from the repository's Scripture chronology corpus, with
its sources named there. The record holds no date for the Sermon on the
Mount or for the occasion of Psalm~32.

\textbf{Rights.} The Latin is published on 17 U.S.C. 103(b) together with its
public-domain witnesses, the Mame printing of 1922 and the Pustet printing of
1862, on the pre-1931 line recorded in \texttt{guidance/missals.md}, applied
as of 26~October 2026. The Douay--Rheims and the Lasance English (public
domain in the United States by nonrenewal) are quoted from registered
witnesses; the NPNF, Library of the Fathers, Migne, Venice, Parma and Marshall
editions are nineteenth-century or earlier publications, Fromage's volume was
published in Paris in 1909, and Schuster's in London in 1930. Scripture,
liturgical texts and quoted translations are not relicensed by the
repository's licence.

\textbf{Use.} This is a study aid to the 1962 Missal, not an official
liturgical book, translation or critical edition.
\par}

\section{References}
\label{sec:references}

{\small
\begin{itemize}[leftmargin=1.2em,itemsep=.25em]
\item \work{Missale Romanum ex decreto Sacrosancti Concilii Tridentini
restitutum, Summorum Pontificum cura recognitum}, editio typica (Typis
Polyglottis Vaticanis, 1962), CMAA facsimile: \work{Rubricae generales} and
\work{Rubricae generales Missalis romani}; Proprium de Tempore, Vigil of the
Nativity, p.~16, and Twenty-third Sunday after Pentecost, pp.~416--417,
nos.~1720, 1726, 1728; \work{Praefationes}, p.~293, no.~1082, and p.~297,
no.~1086; Sanctorale, pp.~702--703, nos.~3974--3983.
\item \work{Missale Romanum}, editio IV iuxta typicam (Tours: Mame, 1922),
pp.~748--750.
\item \work{Missale Romanum} (Ratisbon: Pustet, 1862), pp.~357--358, and the
feast's Offertory (text layer).
\item \work{Missale Romanum} (Venice, 1570), the Offertory of All Saints and
the chants of Ss Marcellinus, Peter and Erasmus and of Ss Felix and Adauctus;
\work{Missale Romanum}, editio typica (Rome, 1604), and \work{Missale Romanum}
(Tournai: Desclée, 1889), the Offertory of All Saints (text layers).
\item Lasance, F.~X., \work{The New Roman Missal in Latin and English},
revised edition (New York: Benziger Brothers, 1945), pp.~733, 735--736
(Twenty-third Sunday after Pentecost) and 1265--1268 (All Saints).
\item \work{The Holy Bible Translated from the Latin Vulgate} (Douay--Rheims,
Challoner revision), Project Gutenberg text: Gen 49:17; 1~Kings 21:10--15; Ps
32:1, 12; 33:1, 8--11; 118:109, 120; Wis 2:10, 20; 3:1--4, 7; Ecclus 27:6; Is
33:6; Ezech 9:4; Amos 8:11; Mt 4:23--25; 5:1--12; 11:25--30; Lk 1:78; 17:21;
Jn 6:48; Rom 8:24; 9:4; 1~Cor 4:11; 2~Cor 1:7; 11:25--27; Gal 3:29; Phil
4:11--12; Apoc 1:9; 7:1--14.
\item \work{Biblia Sacra Vulgatae Editionis} (Clementine Vulgate), for the
comparisons with the Missal's forms.
\item Augustine, \work{Enarrationes in Psalmos} (\work{Expositions on the Book
of Psalms}), NPNF, first series, vol.~8 (CCEL text): on Pss 32 (NPNF ``Psalm
XXXIII'') 1, 12; 33 (NPNF ``Psalm XXXIV'') 13--14; 118 (NPNF ``Psalm CXIX''),
on v.~109.
\item Augustine, \work{De doctrina christiana} III.35.50--51, NPNF, first
series, vol.~2 (CCEL text); \work{Letter XCIII} (to Vincentius), \S30, NPNF,
first series, vol.~1 (CCEL text).
\item Augustine, \work{De sermone Domini in monte} I.1--4, and Sermons 69
(NPNF ``Sermon XIX''), 1--3, and LXI (Benedictine CXI), 1, NPNF, first
series, vol.~6 (CCEL text).
\item Basil the Great, \work{Homilia in Psalmum XXXIII}, on vv.~10--11, PG 29,
cols.~365--368 (Latin version; text layer).
\item Bede, \work{Explanatio Apocalypsis}, book I, on 7:2--17, in E.~Marshall's
English, \work{The Explanation of the Apocalypse} (Oxford and London: James
Parker, 1878).
\item Bellarmine, Robert, \work{Explanatio in Psalmos}, in John O'Sullivan's
English, \work{A Commentary on the Book of Psalms}, on Pss 32 and 33.
\item Chrysostom, John, \work{Homilies on Matthew} 15, 28 and 38, NPNF, first
series, vol.~10 (CCEL text).
\item Cyprian of Carthage, \work{Exhortation to Martyrdom, addressed to
Fortunatus} \S11, ANF, vol.~5 (CCEL text).
\item Durandus, William, \work{Rationale divinorum officiorum} VII.34 and
VII.47, in the 1568 printing with Beleth (Internet Archive OCR).
\item Fromage, Lucien, continuing Prosper Guéranger, \work{L'Année
liturgique}, vol.~15 (Paris: Oudin, 1909), \latin{La Toussaint}, \latin{A la
Messe}, p.~94, citing Gregory the Great, \work{Homiliae in Evangelia} 8.
\item Gregory the Great, \work{Moralia in Iob} XII.v.6, \work{Morals on the
Book of Job}, Library of the Fathers (Oxford, 1844--1850).
\item Hilary of Poitiers, \work{In Evangelium Matthaei commentarius}, cap.~IV,
PL 9, cols.~934--935 (Internet Archive text).
\item Rabanus Maurus, \work{Expositio in Matthaeum}, on Mt 5--6, PL 107
(Internet Archive text).
\item Schuster, Ildefonso, \work{The Sacramentary (Liber Sacramentorum)},
vol.~V (London: Burns Oates \& Washbourne, 1930), pp.~208--212 (Internet
Archive OCR).
\item Thomas Aquinas, \work{Super Evangelium S.~Matthaei lectura}, cap.~5,
lect.~1, in \work{Opera omnia}, vol.~3 (Venice: Bettinelli, 1745);
\work{Super Psalmos}, on Ps 33, nn.~10--11, in \work{Opera omnia}, vol.~14
(Parma, 1863) (Internet Archive OCR).
\item Victorinus of Pettau, \work{Commentary on the Apocalypse of the Blessed
John}, on 7:2, 7:9 and 12:6, with the editors' elucidations, ANF, vol.~7
(CCEL text).
\item \work{The Catholic Encyclopedia} (New York: Robert Appleton, 1907--1912):
Corbett, ``David'' (vol.~4); Durand, ``The New Testament'' (vol.~14); Gigot,
``Book of Wisdom'' (vol.~15); Jacquier, ``Gospel of St.~Matthew'' (vol.~10);
Van den Biesen, ``Apocalypse'' (vol.~1).
\item Haydock, George Leo, \work{The Holy Bible} with Haydock's commentary
(Loreto reprint, 2014), the chronological note at 1~Kings 21 reporting
Ussher.
\item \work{New American Bible, Revised Edition}, introduction to the Psalms
(USCCB website), summarized.
\end{itemize}
\par}
